\mhdl{} provides control syntax start with keyword ``\texttt{metahdl}'',
which interfaces with \mhdlc{} and controls the runtime behavior of compiler.
Designers' controls are passed to compiler via variable assignments embedded
in RTL code, this variable settings are also preceded by keyword \texttt{metahdl}.
Boolean variables inside compiler are set via $+$/$-$ preceded by variable name,
where $+$ means ``enable'' and $-$ means ``disable''.
There are two special form of control syntax: exit syntax, and echo syntax.
The former is used to command compiler exit when the statement is encountered.
The latter is used to print messages on \texttt{stderr}. They are usually used
with preprocessor to guarantee correct configuration settings.
Refer to \autoref{sec:bnf} for detailed formal syntax.

Working scope of all variables can be \emph{Modular}\index{modular variable} or \emph{Effective}\index{effective variable}.
Modular variables (MVAR)\index{MVAR} take effect on entire module and are used when parsing is finished.
Designers can set MVAR anywhere in source code and get the same
effect. If an MVAR is assigned multiple times, last assignment wins.
MVAR can have different values in different files, so file is
the minimum granularity of MVAR.

Effective variables (EVAR)\index{EVAR} take effect from the point the variable is assigned
and are used \emph{during} parsing. Designers can set different values for same EVAR in different
sections of source code, and make compiler treat sections differently. So the minimum granularity
of EVAR is section divided by EVAR assignments.

\section{Variable List}
Following is the complete list of all compiler variables can be assigned by user control
syntax, variable type (boolean or string) and variable scope (MVAR or EVAR) are listed
with variable name.

\begin{description}[style=nextline,leftmargin=2em]

\item[\texttt{modname}\hfill MVAR\,/\,string]
Default: Base file name.\\
Set the generated module name. Often used with preprocessor to
distinguish modules with different configurations.

\item[\texttt{outfile}\hfill MVAR\,/\,string]
Default: Base file name.\\
Set the generated \sv{} file base name. Often used with preprocessor to
distinguish module definition files with different configurations.

\item[\texttt{portchk}\hfill MVAR\,/\,boolean]
Default: \texttt{false}.\\
Enable/Disable port validation for module.

\item[\texttt{hierachydepth}\hfill MVAR\,/\,positive int]
Default: \texttt{300}.\\
Maximum level of module instantiation.

\item[\texttt{clock}\hfill EVAR\,/\,string]
Default: \texttt{clk}.\\
Default clock name used for \texttt{ff\_block} and \texttt{fsm\_block}.

\item[\texttt{reset}\hfill EVAR\,/\,string]
Default: \texttt{rst\_n}.\\
Default reset name used for \texttt{ff\_block} and \texttt{fsm\_block}.

\item[\texttt{multidriverchk}\hfill MVAR\,/\,boolean]
Default: \texttt{true}.\\
Enable/Disable multiple driver checking for module.

\item[\texttt{relexedfsm}\hfill EVAR\,/\,boolean]
Default: \texttt{true}.\\
Set severity of connectivity/reachability error checked in FSM.
If it is true, relaxed FSM programming mode is enabled,
all dead states or unreachable states are acceptable, compiler
only reports warning when such states are encountered, and continues processing.
If it is false, FSM programming is in strict mode, any dead state or
unreachable state is considered to be fatal error, compiler will report error
and stop processing if such state is checked.

\item[\texttt{exitonwarning}\hfill EVAR\,/\,boolean]
Default: \texttt{false}.\\
Set severity of \emph{normal parsing warning}, such as width mismatch.
If it is true, compiler exits on any warning.

\item[\texttt{exitonlintwarning}\hfill MVAR\,/\,boolean]
Default: \texttt{false}.\\
Warning from port validation, multiple driver checking are categorized
as lint warning. If this variable is set to true, compiler will exit
on any lint warning.

\end{description}

\section{Example}
\autoref{lst:top wrapper} demonstrates user control syntax with
code configuration.

\begin{lstlisting}[caption={Demo Top Wrapper}, label={lst:top wrapper}]
metahdl + portchk;

/*-\pp{if}-*/ WIDTH > 64
metahdl ``width can not exceed 64!'';
metahdl exit;
/*-\pp{endif}-*/

assign data[/*-\pp{WIDTH}-*/-1:0] = /*-\pp{WIDTH}-*/'d0;

/*-\pp{ifdef}-*/ FPGA
/*-\pp{define}-*/ target fpga
/*-\pp{else}-*/
/*-\pp{define}-*/ target asic
/*-\pp{endif}-*/

metahdl modname = top_/*-\`{}\textcolor{purple}{target}-*/;
metahdl outfile = top_/*-\`{}\textcolor{purple}{target}-*/;

metahdl + exitonwarning;

metahdl clock = clk_125M;
metahdl reset = pclk_rst_n;
ff;
  a_ff, a, 1'b0;
  b_ff[1:0], b, 1'b0;
endff
metahdl - exitonwarning;

metahdl clock = clk_250M;
metahdl reset = dclk_rst_n;
ff;
  c_ff, c, 1'b0;
  d_ff, d, 1'b0;
endff

ff;
  e_ff, e, 1'b0;
  g_ff, g, 1'b0;
endff
\end{lstlisting}

In this example, port validation is enabled for the module.
The \texttt{\`{}if}/\texttt{\`endif} block checks the value of macro \texttt{WIDTH}
and forces compilation exit upon illegal values.
The \texttt{\`{}ifdef}/\texttt{\`else}/\texttt{\`endif} block defines the target device
type, and \texttt{modname}/\texttt{outfile} variables alter \sv{} module name and
output file name accordingly.

The \texttt{exitonwarning} variable is enabled only for the first \texttt{ff} block,
so width mismatch between \verb|b_ff| and \verb|b| is treated as a fatal error.
Different clock and reset names (\texttt{clk\_125M}/\texttt{pclk\_rst\_n} vs.\
\texttt{clk\_250M}/\texttt{dclk\_rst\_n}) are used for the two \texttt{ff} sections,
demonstrating how EVAR scope works.
